\subsection{PRODUCT SPACES}

DEFINITION 1. Given a family \((\mathbf{X}_{t})_{t\in I}\) of topological spaces, the product space of this family is the product set \(\mathbf{X} = \prod_{t\in I}\mathbf{X}_{t}\) with the topology which is the product of the topologies of the \(\mathbf{X}_{t}\) (§ 2, no. 3, Example 3). The spaces \(\mathbf{X}_{t}(t\in I)\) are called the factors of \(\mathbf{X}\) .

By virtue of § 2, no. 3, Proposition 4, the product topology on X has as a base the set B of finite intersections of sets of the form \(\overline{\mathrm{pr}}_{t}^{1}(\mathbf{U}_{t})\) , where \(U_{t}\) is open in \(X_{t}\) ; these sets are products \(\prod_{i\in I}A_{i}\) , where \(A_{i}\) is open in \(X_{t}\) for each \(i\in I\) and \(A_{i}=X_{t}\) for all but a finite number of indices. These sets will be called elementary sets.

If \(\mathfrak{B}_i\) is a base of the topology of \(X_i\) (for each \(i \in I\) ), it is clear that the elementary sets \(\prod_{i \in I} A_i\) such that \(A_i \in \mathfrak{B}_i\) for each index \(i\) such that \(A_i \neq X_i\) form another base of the product topology. The elementary sets of this type which contain a given point \(x \in X\) thus form a fundamental system of neighbourhoods of \(x\) ( \(\S 1\) , no. 3, Proposition 3).

If I is a finite set, the construction of the product topology from the topologies of the factors \(X_{i}\) is simpler: the elementary sets are just products \(\prod_{i\in I}A_{i}\) , where \(A_{i}\) is any open subset of \(X_{i}\) , for each \(i\in I\) (cf.~Exercise 9).

Examples. * The product \(\mathbf{R}^n\) of \(n\) spaces identical with the real line \(\mathbf{R}\) is called real number space of \(n\) dimensions; \(\mathbf{R}^2\) is also called the real plane (cf.~Chapter VI, § 1). Likewise, starting from the rational line \(\mathbf{Q}\) , we define the rational number space of \(n\) dimensions \(\mathbf{Q}^n\) (rational plane for \(n = 2\) ).

The topology of the space \(\mathbf{R}^n\) has as a base the set of all products of \(n\) open intervals in \(\mathbf{R}\) , which are called open boxes of \(n\) dimensions. The open boxes which contain a point \(x \in \mathbf{R}^n\) form a fundamental system of neighbourhoods of this point. Likewise the products of \(n\) closed intervals in \(\mathbf{R}\) are called closed boxes of \(n\) dimensions. The closed boxes to which \(x\) is interior also form a fundamental system of neighbourhoods of \(x\) . There are analogous results for \(\mathbf{Q}^n\) . *

PROPOSITION 1. Let \(f = (f_{i})\) be a mapping of a topological space \(Y\) into a product space \(X = \prod_{i \in I} X_{i}\) . Then \(f\) is continuous at a point \(a \in Y\) if and only if \(f_{i}\) is continuous at a for each \(i \in I\) .

Since \(f_{t} = \mathrm{pr}_{t} \circ f\) , this is just a particular case of Proposition 4 of §2, no. 3.

COROLLARY 1. Let \((\mathbf{X}_t)_{i \in \mathbf{I}}\) , \((\mathbf{Y}_t)_{i \in \mathbf{I}}\) be two families of topological spaces with the same set of indices. For each \(i \in \mathbf{I}\) , let \(f_i\) be a mapping of \(\mathbf{X}_t\) into \(\mathbf{Y}_t\) . In order that the product mapping \(f: (x_i) \to (f_i(x_i))\) of

\[
\prod_ {i \in I} X _ {i} \quad i n t o \quad \prod_ {i \in I} Y _ {i}
\]

should be continuous at a point \(a = (a_{i})\) it is necessary and sufficient that \(f_{i}\) is continuous at \(a_{i}\) for each \(\mathfrak{t} \in \mathbf{I}\) .

\(f\) can be written as \(x \to (f_{t}(\mathrm{pr}_{t}(x))\) , so that by Proposition 1 the condition is sufficient. Conversely, for each \(x \in I\) let \(g_{x}\) be the mapping of \(X_{x}\) into \(\prod_{i \in I} X_{i}\) such that \(\mathrm{pr}_{x}(g_{x}(x_{x})) = x_{x}\) and \(\mathrm{pr}_{t}(g_{x}(x_{x})) = a_{t}\) whenever \(i \neq x\) ; then \(g_{x}\) is continuous at the point \(a_{x}\) , by Proposition 1. Since \(f_{x} = \mathrm{pr}_{x} \circ f \circ g_{x}\) it follows that if \(f\) is continuous at \(a\) , then \(f_{x}\) is continuous at \(a_{x}\) .

COROLLARY 2. Let X, Y be two topological spaces. In order that a mapping \(f: \mathbf{X} \to \mathbf{Y}\) should be continuous it is necessary and sufficient that the mapping \(g: x \to (x, f(x))\) is a homeomorphism of X onto the graph G of \(f\) (considered as a subspace of the product space \(\mathbf{X} \times \mathbf{Y}\) ).

Since \(f = \mathrm{pr}_2 \circ g\) , the condition is sufficient. It is also necessary, for if \(f\) is continuous, then \(g\) is bijective and continuous (Proposition 1) and the inverse of \(g\) is the restriction of \(\mathrm{pr}_1\) to \(G\) , which is continuous (cf.~Set Theory, Chapter IV, § 2, no. 4, criterion CST 17).

PROPOSITION 2 (Associativity of topological products). Let \((\mathbf{X}_{\iota})_{\iota \in \mathbf{I}}\) be a family of topological spaces, \((\mathbf{J}_{\kappa})_{\kappa \in \mathbf{K}}\) a partition of the set I, and for each \(\kappa \in \mathbf{K}\) let \(\mathbf{X}_{\kappa}^{\prime} = \prod_{i\in \mathbf{J}_{\kappa}}\mathbf{X}_{\iota}\) be the product of the spaces \(\mathbf{X}_{\iota}\) for \(\iota \in \mathbf{J}_{\mathbf{K}}\) . Then the canonical mapping of the product space
% semantic-fragments: legacy-input-seam
\relax{}
\[\prod_{i \in I} X_i\ \text{ onto the product space }\ \prod_{x \in K} X_x'\]

is a homeomorphism.

This is a particular case of transitivity of initial topologies (§ 2, no. 3, Proposition 5; cf.~Set Theory, Chapter IV, § 2, no. 4, criterion CST 13).

Generally we identify the product spaces \(\prod_{i\in I}X_i\) and \(\prod_{x\in X}X_x'\) by means of the canonical mapping.

COROLLARY. Let \(\sigma\) be a permutation of the set I. Then the mapping \((x_{i})\to (x_{\sigma(i)})\) is a homeomorphism of

\[
\prod_ {i \in I} X _ {i} \quad onto \quad \prod_ {i \in I} X _ {\sigma (i)}.
\]

Take \(\mathbf{K} = \mathbf{I}\) and \(\mathbf{J}_t = \{\sigma(t)\}\) for each \(t \in I\) in Proposition 2.

PROPOSITION 3. Let X be a set, \((\mathbf{Y}_{t})_{t\in I}\) a family of topological spaces, and for each \(t\in I\) let \(f_{t}\) be a mapping of X into \(Y_{t}\) . Let f be the mapping \(x\to(f_{t}(x))\) of X into \(Y=\prod_{t\in I}Y_{t}\) , and let G be the coarsest topology on X for which the mappings \(f_{t}\) are continuous. Then G is the inverse image under f of the topology induced on \(f(\mathbf{X})\) by the product topology on Y.

This is another particular case of transitivity of initial topologies (§ 2, no. 3, Proposition 5; cf.~Set Theory, Chapter IV, § 2, no. 4, criterion CST 15).

COROLLARY. For each \(t \in I\) let \(A_{t}\) be a subspace of \(Y_{t}\) . Then the topology induced on \(A = \prod_{t \in I} A_{t}\) by the product topology on \(\prod_{t \in I} Y_{t}\) is the product of the topologies of the subspaces \(A_{t}\) .

Let \(j_{i}\) be the canonical injection \(A_{i} \to Y_{i}\) , and apply Proposition 3 to the mappings \(f_{i} = j_{i} \circ \operatorname{pr}_{i}\) (cf.~Set Theory, Chapter IV, § 2, no. 4, criterion GST 14).

